% ---------------------------------------------------------------------------
% Resultados da RSL 2 (assistentes e agentes de IA no desenvolvimento).
% Base: docs/triagem/ (prisma-etapa2, snowball-leitura, qualidade, sintese).
% Situação descrita: 12/11/2026.
% ---------------------------------------------------------------------------

\section{RSL 2: assistentes e agentes de IA no desenvolvimento de software}
\label{sec:resultados-rsl2}

A leitura do texto completo na etapa 2, o \textit{snowballing}, a avaliação
de qualidade, a extração e a síntese desta revisão foram feitos pela IA como
terceira opinião, com regras aceitas pelo autor (R1). O orientador (R2) ainda
confirma essas regras, e a avaliação independente de R1 e R2, com o kappa de
cada etapa, ainda não foi feita. Os resultados abaixo podem mudar quando isso
acontecer. A unidade da síntese é o estudo, e não o relato: um estudo pode ter
mais de um relato publicado. Os estudos incluídos são citados pelo
identificador entre parênteses, por exemplo (RSL2-0487), e a lista completa
está no Apêndice~\ref{ap:estudos}.
% PENDENTE: R2 confirma as regras R2-A a R2-H (seleção), as do snowballing,
% Q-A a Q-D (qualidade) e S-A a S-E (síntese). A avaliação independente de
% R1 e R2 e o kappa da etapa 2 e da qualidade ainda não foram feitos.

% ---------------------------------------------------------------------------
\subsection{Seleção dos estudos}
\label{sec:resultados-rsl2-selecao}

% Gerado por docs/triagem/sintese/scripts/gera_tex.py. Não editar à mão.
\begin{table}[htbp]
\centering
\caption{Fluxo de seleção da RSL 2: busca nas bases}
\label{tab:prisma-rsl2}
\small
\begin{tabular}{p{2.2cm}p{9.5cm}r}
\hline
\textbf{Fase} & \textbf{Caixa} & \textbf{n} \\
\hline
Identificação & Registros das bases & 8.181 \\
 & Scopus & 3.148 \\
 & Web of Science & 1.914 \\
 & IEEE Xplore (depois do filtro de tipo) & 1.577 \\
 & ACM DL (depois do filtro de tipo) & 1.412 \\
 & Fontes nacionais (SBES e JSERD) & 130 \\
 & Removidos antes da triagem: duplicatas (3.857 pelo script, 1 apagada no Rayyan) & 3.858 \\
Triagem & Registros triados (título, resumo e palavras-chave) & 4.323 \\
 & Excluídos na etapa 1 & 3.208 \\
 & Passaram para a etapa 2 & 1.115 \\
 & Excluídos pelo recorte, pelo resumo e metadados (CI6 583, dos quais 21 pela trilha na Crossref/ACM DL; CI5 280) & 863 \\
 & Estudos secundários, analisados à parte & 63 \\
 & Relatos buscados (texto completo) & 189 \\
 & Relatos não recuperados (CE6: UnB/CAPES e acesso aberto) & 12 \\
 & Relatos lidos (elegibilidade) & 177 \\
 & Excluídos no texto completo (CI5 35, CI2 23, CE4 9, CI6 8, CI1 5, CE2 2, CE1 1, CE3 1, CE5 1, CE7 1, CI3 1) & 87 \\
Inclusão & Estudos primários incluídos & 90 \\
\hline
\end{tabular}
\fonte{Elaborada pelo autor, no modelo de fluxograma do PRISMA 2020 \cite{page2021prisma}.}
\end{table}
\begin{table}[htbp]
\centering
\caption{Fluxo de seleção da RSL 2: \textit{snowballing}}
\label{tab:snowball-rsl2}
\small
\begin{tabular}{p{11.7cm}r}
\hline
\textbf{Caixa} & \textbf{n} \\
\hline
Registros novos identificados (referências e citações) & 8.186 \\
Excluídos pelos metadados & 5.418 \\
Triados por título e resumo & 2.768 \\
Excluídos na triagem & 2.501 \\
Estudos secundários (analisados à parte) & 58 \\
Relatos buscados & 209 \\
Relatos não recuperados & 35 \\
Relatos avaliados no texto completo & 174 \\
Excluídos no texto completo & 134 \\
\textbf{Estudos incluídos pelo \textit{snowballing}} & \textbf{40} \\
\hline
\end{tabular}
\fonte{Elaborada pelo autor, na coluna ``outros métodos do modelo do PRISMA 2020 \cite{page2021prisma}.}
\end{table}


A Tabela~\ref{tab:prisma-rsl2} mostra o fluxo da busca nas bases. As bases
devolveram 8.181 registros. Depois da remoção de 3.858 duplicatas, 4.323
registros foram triados por título, resumo e palavras-chave, e 3.208 foram
excluídos na etapa 1. Dos 1.115 que passaram, 863 saíram pelo recorte da
etapa 2, decidido pelo resumo e pelos metadados: só ficaram estudos com
evidência de desenvolvedores profissionais (CI5) publicados em periódico ou
em conferência A* ou A no CORE 2023 (CI6)
(Seção~\ref{sec:metodologia-selecao}). Outros 63 registros eram estudos
secundários e foram separados para análise à parte. Foram buscados 189
relatos, 12 não foram recuperados, 177 foram lidos e 87 foram excluídos no
texto completo, o que deixou 90 relatos incluídos pela busca nas bases.

O \textit{snowballing} (Tabela~\ref{tab:snowball-rsl2}) identificou 8.186
registros novos e parou na terceira iteração, a primeira sem inclusão nova.
Dele vieram mais 40 relatos, de modo que a revisão tem 130 relatos incluídos.
Dois deles, SB2-1833 e SB2-1960, têm o mesmo autor, o mesmo caso e os mesmos
resultados, publicados em dois periódicos. Por isso contam como um só estudo,
com SB2-1960 proposto como relato principal, porque é o que traz a seção
de limitações. A síntese trata, portanto, de 129 estudos. A avaliação de
qualidade, que é feita por relato, continua com 130.
% PENDENTE: R1 e R2 decidem o par SB2-1833/SB2-1960. Se o par for mantido, a
% coluna do snowballing tem 40 relatos e 39 estudos; o fluxograma precisa
% indicar isso.
% PENDENTE: RSL2-4104 é extensão em periódico de um artigo dos mesmos autores
% no ICSE 2025, com a mesma amostra. Uma busca pelos títulos não achou o
% artigo do ICSE entre os incluídos; R1 confirma pelos autores que não há
% dupla contagem.

Duas limitações da seleção afetam a leitura dos resultados. A primeira são os
relatos não recuperados: 12 na busca nas bases e 35 no
\textit{snowballing}, que não puderam ser lidos \cite[item 23c]{page2021ee}.
A segunda são os estudos secundários (63 na busca e 58 no
\textit{snowballing}), cujo texto completo ainda não foi lido e que serão
analisados à parte, como prevê o protocolo.

% ---------------------------------------------------------------------------
\subsection{Características dos estudos}
\label{sec:resultados-rsl2-mapeamento}

% Gerado por docs/triagem/sintese/scripts/gera_tex.py. Não editar à mão.
\begin{table}[htbp]
\centering
\caption{RSL 2: ano de publicação dos estudos incluídos}
\label{tab:map-rsl2-anodepublica}
\small
\begin{tabular}{p{8cm}rr}
\hline
\textbf{Categoria} & \textbf{n} & \textbf{\%} \\
\hline
2022 & 1 & 1 \\
2023 & 3 & 2 \\
2024 & 23 & 18 \\
2025 & 32 & 25 \\
2026 & 66 & 51 \\
2027 & 4 & 3 \\
\hline
\end{tabular}
\fonte{Elaborada pelo autor.}
\end{table}
\begin{table}[htbp]
\centering
\caption{RSL 2: tipo de estudo dos estudos incluídos}
\label{tab:map-rsl2-tipodeestudo}
\small
\begin{tabular}{p{8cm}rr}
\hline
\textbf{Categoria} & \textbf{n} & \textbf{\%} \\
\hline
métodos mistos & 40 & 31 \\
entrevistas/qualitativo & 23 & 18 \\
experimento controlado & 23 & 18 \\
survey & 16 & 12 \\
estudo de campo/caso & 15 & 12 \\
outro & 7 & 5 \\
quase-experimento & 5 & 4 \\
\hline
\end{tabular}
\fonte{Elaborada pelo autor.}
\end{table}
\begin{table}[htbp]
\centering
\caption{RSL 2: contexto dos estudos incluídos}
\label{tab:map-rsl2-contextoqp}
\small
\begin{tabular}{p{8cm}rr}
\hline
\textbf{Categoria} & \textbf{n} & \textbf{\%} \\
\hline
indústria & 65 & 50 \\
misto & 42 & 33 \\
academia & 15 & 12 \\
código aberto & 7 & 5 \\
\hline
\end{tabular}
\fonte{Elaborada pelo autor.}
\end{table}
\begin{table}[htbp]
\centering
\caption{RSL 2: nível de autonomia dos estudos incluídos}
\label{tab:map-rsl2-nveldeautono}
\small
\begin{tabular}{p{8cm}rr}
\hline
\textbf{Categoria} & \textbf{n} & \textbf{\%} \\
\hline
misto & 64 & 50 \\
assistente conversacional & 39 & 30 \\
agente autônomo & 14 & 11 \\
completar código & 10 & 8 \\
outro/não relatado & 2 & 2 \\
\hline
\end{tabular}
\fonte{Elaborada pelo autor.}
\end{table}


O corpus é recente (Tabela~\ref{tab:map-rsl2-anodepublica}): 66 dos 129
estudos (51\%) são de 2026, 32 de 2025 e 23 de 2024. Quatro estudos têm o ano
2027 nos metadados. Há 79 estudos em periódicos e 50 em conferências. O país
principal (o primeiro citado) mais frequente é os EUA (38 estudos), e cinco
estudos são do Brasil.
% PENDENTE: conferir se o ano 2027 de RSL2-0120, RSL2-0201, RSL2-0246 e
% RSL2-0322 (OpenAlex ou base) é o ano do volume.

Os desenhos são variados (Tabela~\ref{tab:map-rsl2-tipodeestudo}). Métodos
mistos são o tipo mais comum (40 estudos), seguidos de entrevistas e estudos
qualitativos (23) e de experimentos controlados (23). Há ainda 16
\textit{surveys}, 15 estudos de campo ou de caso, 5 quase-experimentos e 7
estudos de outros tipos. Metade dos estudos vem da indústria (65), 42 têm
contexto misto, 15 são da academia e 7 de projetos de código aberto
(Tabela~\ref{tab:map-rsl2-contextoqp}).

O nível de autonomia responde à primeira parte da QP6
(Tabela~\ref{tab:map-rsl2-nveldeautono}). Metade dos estudos (64) trata de
mais de um nível. Assistentes conversacionais são o foco de 39 estudos,
agentes autônomos de 14 e o completar código de 10. A leitura B4
(Tabela~\ref{tab:temas-B4}, TA6) mostra que a evidência sobre ferramentas é
heterogênea e datada. Copilot e ChatGPT predominam, muitos estudos usam
ferramenta própria, e a versão ou a data de uso muitas vezes não é relatada
(RSL2-0120, RSL2-2903, SB2-1143). Os agentes autônomos aparecem sobretudo em
ferramentas internas de grandes empresas, em reparo, testes e revisão
(RSL2-0150, SB2-0162, SB2-2613, RSL2-0322), sempre com decisão humana antes
da integração. Nenhum estudo avalia um agente que integre código sem esse
portão, de modo que a autonomia plena não foi estudada. O caso
mais próximo de um projeto quase \textit{greenfield} com agente é RSL2-0487,
com Claude Code, decomposição em itens atômicos e revisão cruzada entre
modelos. Com desenvolvedores cegos, a preferência migrou do modo
\textit{Agent} para o \textit{Ask} com confirmação (RSL2-0750).

% ---------------------------------------------------------------------------
\subsection{Qualidade dos estudos}
\label{sec:resultados-rsl2-qualidade}

As notas da RSL 2 estão nas Tabelas~\ref{tab:qualidade-resumo}
e~\ref{tab:qualidade-itens}. A numeração dos critérios da RSL 2 difere da RSL
1, e a correspondência está na Tabela~\ref{tab:qualidade}. Os 130 relatos
têm nota média de 9,46 e mediana 10, numa escala de 0 a 11, com notas entre 4
e 11. Há 79 relatos com nota 10 ou 11, 43 com nota de 7 a 9 e 8 com nota de 4
a 6. As notas mais baixas (4 e 5) são de SB2-1833, SB2-1960, RSL2-1040,
RSL2-1982, RSL2-2903 e RSL2-4029. Como prevê a
Seção~\ref{sec:metodologia-qualidade}, a nota pondera a síntese e não exclui
estudos.

As notas se concentram no topo da escala, um efeito teto. O recorte da etapa
2 já selecionou estudos bem relatados, com desenvolvedores profissionais,
publicados em periódicos ou em conferências A* ou A. Por isso a nota diferencia pouco os estudos
desta revisão, e a análise de sensibilidade só com notas $\geq$ 8
(Seção~\ref{sec:resultados-rsl2-direcao}) quase não muda a contagem. Os três
critérios com menos respostas ``sim'' indicam onde estão as fraquezas do
corpus. O AQ5 (grupo de controle) tem 67 de 130 (52\%), porque muitos estudos
são \textit{surveys} ou estudos qualitativos. O AQ8 (relação entre
pesquisador e participantes) tem 69 de 130 (53\%). O AQ11, acréscimo desta
revisão, tem 94 de 130 (72\%): cerca de um quarto dos relatos não informa o
modelo e a versão (ou a data) da ferramenta usada, o que dificulta comparar
e replicar os resultados.
% PENDENTE: o código C38 da leitura B4 ("versão ou data não relatada", que
% inclui a omissão parcial) tem uma lista de estudos maior que o número de
% "não" no AQ11. R1 e R2 conferem se a diferença vem só da omissão parcial.

% ---------------------------------------------------------------------------
\subsection{Direção do efeito}
\label{sec:resultados-rsl2-direcao}

% Gerado por docs/triagem/sintese/scripts/gera_tex.py. Não editar à mão.
\begin{table}[htbp]
\centering
\caption{RSL 2: direção do efeito relatado (entre parênteses, só estudos com nota $\geq$ 8)}
\label{tab:direcao-rsl2}
\small
\begin{tabular}{llrrrrrr}
\hline
\textbf{Dimensão} & \textbf{QP} & \textbf{Positivo} & \textbf{Negativo} & \textbf{Misto} & \textbf{Nulo} & \textbf{Não medido} & \textbf{Mediram} \\
\hline
produtividade & QP1 & 51 (46) & 4 (4) & 29 (27) & 5 (5) & 40 (37) & 89 \\
qualidade & QP2 & 22 (20) & 12 (11) & 56 (51) & 4 (4) & 35 (33) & 94 \\
segurança & QP3 & 3 (3) & 14 (12) & 4 (4) & 2 (2) & 106 (98) & 23 \\
custo & QP4 & 3 (3) & 2 (2) & 0 (0) & 2 (2) & 122 (112) & 7 \\
\hline
\end{tabular}
\fonte{Elaborada pelo autor. Contagem pela direção do efeito \cite[quadro 1]{page2021ee}.}
\end{table}
\begin{table}[htbp]
\centering
\caption{RSL 2: direção do efeito por tipo de medida}
\label{tab:direcao-medida-rsl2}
\small
\begin{tabular}{llrrrr}
\hline
\textbf{Dimensão} & \textbf{Medida} & \textbf{Positivo} & \textbf{Negativo} & \textbf{Misto} & \textbf{Nulo} \\
\hline
produtividade & medido & 13 & 3 & 5 & 5 \\
produtividade & percebido & 18 & 0 & 11 & 0 \\
produtividade & ambos & 20 & 1 & 13 & 0 \\
qualidade & medido & 11 & 6 & 14 & 3 \\
qualidade & percebido & 3 & 5 & 27 & 0 \\
qualidade & ambos & 8 & 1 & 15 & 1 \\
\hline
\end{tabular}
\fonte{Elaborada pelo autor. Só estudos que mediram a dimensão. ``Ambos: métrica objetiva e percepção no mesmo estudo.}
\end{table}


A Tabela~\ref{tab:direcao-rsl2} conta os estudos pela direção do efeito
relatado em cada dimensão \cite[quadro 1]{page2021ee}. A contagem dá a
direção, não o tamanho do efeito, e um estudo ``positivo'' pode ter um efeito
pequeno. Em produtividade (QP1), 89 estudos mediram o efeito: 51 têm direção
positiva, 29 mista, 4 negativa e 5 nula. Em qualidade (QP2), 94 mediram, e a
direção mista predomina (56), contra 22 positivos, 12 negativos e 4 nulos.
Segurança (QP3) foi medida em só 23 estudos, dos quais 14 têm direção
negativa, 4 mista, 3 positiva e 2 nula. Os três positivos (RSL2-0383,
RSL2-0487 e SB2-0308) avaliam ferramentas de verificação ou
\textit{pipelines} com análise estática, e não o uso da IA sozinho. Custo
(QP4) foi medido em 7 estudos: 3 positivos (RSL2-0378, RSL2-0487 e
SB2-2374), 2 negativos e 2 nulos. Os números entre parênteses repetem a
contagem só com os estudos de nota $\geq$ 8. Por causa do efeito teto, o
quadro não muda.
% PENDENTE: o corte da análise de sensibilidade (nota >= 8) é decisão do
% autor; R2 confirma.

A Tabela~\ref{tab:direcao-medida-rsl2} separa produtividade e qualidade pelo
tipo de medida: métrica objetiva (medido), autorrelato (percebido) ou as duas
no mesmo estudo (ambos). As três linhas de cada dimensão somam os estudos que
mediram. Em produtividade, quando o efeito é só percebido, há 18 estudos
positivos, 11 mistos e nenhum negativo ou nulo. Quando é medido, há 13
positivos, 5 mistos, 3 negativos e 5 nulos. Em qualidade, 27 dos 35 estudos
com efeito só percebido têm direção mista, enquanto os medidos se dividem em
11 positivos, 6 negativos, 14 mistos e 3 nulos. Nas duas dimensões, os
resultados negativos e nulos aparecem sobretudo quando há medida objetiva, e
a direção é mais favorável quando o efeito é percebido. Classificar a medida
pelo desenho do estudo leva ao mesmo padrão. Esse contraste reaparece nos temas de produtividade (TA1 de B1) e
de qualidade (TA4 de B2).
% PENDENTE: a classificação de cada estudo como medido, percebido ou ambos é
% decisão do autor (leituras B1 e B2, com trecho conferido); R1 e R2 conferem.

% ---------------------------------------------------------------------------
\subsection{QP1: produtividade}
\label{sec:resultados-rsl2-qp1}

% Gerado por docs/triagem/sintese/scripts/gera_tex.py. Não editar à mão.
\begin{table}[htbp]
\centering
\caption{Temas analíticos: RSL 2, QP1}
\label{tab:temas-B1}
\small
\begin{tabular}{lp{11cm}r}
\hline
\textbf{Tema} & \textbf{Descrição} & \textbf{Estudos} \\
\hline
TA1 & O ganho existe, mas é condicional e menor do que se percebe & 88 \\
TA2 & O ganho depende de quem usa, de qual tarefa e de como se interage & 55 \\
TA3 & Custos ocultos consomem parte do ganho & 61 \\
TA4 & As métricas são heterogêneas e muitas medem a ferramenta, não o resultado & 96 \\
\hline
\end{tabular}
\fonte{Elaborada pelo autor. ``Estudos: número de estudos distintos nos códigos do tema.}
\end{table}


A síntese da QP1 seguiu as três etapas de \citeonline{thomas2008}:
códigos, temas descritivos e temas analíticos. A Tabela~\ref{tab:temas-B1}
dá, para cada tema analítico, o número de estudos nos códigos do tema. Esse
número mede a amplitude do tema, e não quantos estudos sustentam cada
afirmação, que vem dos estudos citados no texto.

O primeiro tema (TA1) é que o ganho existe, mas é condicional e menor do que
se percebe. Medidas objetivas mostram menos tempo em tarefas curtas e bem
definidas, como 22\% em RSL2-0772 e $-$63\% na tarefa de OAuth de SB2-2049.
Desenhos controlados também acham efeito nulo (fase 2 de RSL2-0393, RSL2-0559,
RSL2-4042) ou negativo (RSL2-0180, RSL2-1328). Em campo, as horas semanais
caíram sem mudar o número de PRs nem o tempo de ciclo (RSL2-0205). Os valores
mais altos, como 25 a 40 vezes em RSL2-0487 e 84\% em SB2-0308, vêm de casos
únicos sem grupo de controle concorrente. A IA acelera tarefas isoladas, mas
a evidência de que acelera a entrega é fraca. O segundo tema (TA2) mostra que
o ganho depende de quem usa, da tarefa e da forma de interagir. Os ganhos se
concentram em tarefas rotineiras, \textit{boilerplate} e tarefas delimitadas
com verificação automática, como testes e correção guiada por análise
estática (SB2-0162, SB2-0308, SB2-1193, SB2-2107). Em tarefas complexas, o
ganho some ou o desenvolvedor abandona a ferramenta (RSL2-1385, RSL2-1181,
RSL2-3990). Interação incremental supera a geração em uma só vez
(\textit{one-shot}) em RSL2-2485, e \textit{prompt} predefinido supera
\textit{prompt} improvisado em SB2-0471.

O terceiro tema (TA3) é que custos ocultos consomem parte do ganho. Verificar,
corrigir e reescrever a saída ocupou 22,4\% do tempo em RSL2-2593, e a
validação e as correções dominaram o esforço em RSL2-0240. Descarte de código
(RSL2-0246), microinterrupções (RSL2-0214) e o custo de escrever
\textit{prompts} (RSL2-3768, SB2-0690) reduzem o saldo. Na revisão de código
assistida, o tempo medido não cai (RSL2-2484, RSL2-4042, RSL2-3693). O quarto tema (TA4) é
que as métricas são heterogêneas e muitas medem a ferramenta, não o
resultado. A taxa de aceitação de sugestões vai de 22\% (SB2-0206) a 84,26\%
(RSL2-2812) e prediz só a produtividade percebida (RSL2-3300). Métricas de
entrega, como PRs por mês e tempo de ciclo, aparecem só em RSL2-0205, e
empresas relatam não ter métricas para quantificar o ganho (SB2-1783).

As contradições confirmam esses temas. A experiência modera o efeito em
direções opostas: novatos ganham mais em RSL2-0120, RSL2-1982 e SB2-1818, mas
não concluem a tarefa em SB2-0690. As duas fases de RSL2-0393 divergem: a
observacional mede 30,7\% menos tempo com IA, e a randomizada é nula. SB2-1818 mede mais tempo com Bard junto com a percepção de mais
rapidez. As lacunas são três. Faltam experimentos que meçam o tempo do desenvolvedor com
agentes autônomos contra um grupo de controle humano (RSL2-0150, RSL2-1320,
SB2-1193, entre outros). Os experimentos usam tarefas de laboratório de 10 a
45 minutos (RSL2-1182, RSL2-2650, RSL2-0214, RSL2-3387, RSL2-1500). E efeitos
de longo prazo medidos são raros (RSL2-0111, RSL2-0205).

% ---------------------------------------------------------------------------
\subsection{QP2: qualidade do software}
\label{sec:resultados-rsl2-qp2}

% Gerado por docs/triagem/sintese/scripts/gera_tex.py. Não editar à mão.
\begin{table}[htbp]
\centering
\caption{Temas analíticos: RSL 2, QP2}
\label{tab:temas-B2}
\small
\begin{tabular}{lp{11cm}r}
\hline
\textbf{Tema} & \textbf{Descrição} & \textbf{Estudos} \\
\hline
TA1 & Não há ganho uniforme de correção: o efeito depende da tarefa, da condição e do modelo & 67 \\
TA2 & Código que passa nos testes ainda pode ser pior de manter & 48 \\
TA3 & A qualidade depende da verificação humana, e é ela que a IA enfraquece & 48 \\
TA4 & A percepção de qualidade não coincide com a medição & 72 \\
TA5 & Contexto, iteração com feedback de execução e escopo pequeno melhoram a saída & 56 \\
TA6 & Documentação gerada ajuda, mas pede revisão de precisão e de verbosidade & 40 \\
\hline
\end{tabular}
\fonte{Elaborada pelo autor. ``Estudos: número de estudos distintos nos códigos do tema.}
\end{table}


A Tabela~\ref{tab:temas-B2} lista os seis temas analíticos da QP2. O primeiro
(TA1) é que não há ganho uniforme de correção. Medições com uma condição sem
IA apontam nos três sentidos: correção maior (RSL2-0772, RSL2-1182,
SB2-2049), menor (RSL2-2576, SB2-1255) e nula (RSL2-1500, RSL2-2650). A
aprovação cai em tarefas difíceis ou pedidos longos (RSL2-2839, RSL2-2068,
SB2-0279) e varia com o modelo e a versão (RSL2-0488, RSL2-3990, SB2-1605).
O segundo (TA2) é que código que passa nos testes ainda pode ser pior de
manter. Análises estáticas medem mais duplicação, complexidade e
recomendações de manutenibilidade no código feito com IA (RSL2-1246,
RSL2-0393, SB2-1605, SB2-2107). Testes gerados têm cobertura alta, mas podem
ser tautológicos, ter asserções fracas ou ser burlados por correções
automáticas (RSL2-0487, RSL2-0378, RSL2-0150, SB2-0162). Segurança e
acessibilidade quase não entram nas avaliações de qualidade (RSL2-1840,
RSL2-0225).

O terceiro tema (TA3) é que a qualidade depende da verificação humana, e é
ela que a IA enfraquece. Os estudos convergem em exigir revisão humana antes
de integrar (RSL2-0761, SB2-1143, RSL2-1385, RSL2-0180). Medições mostram que
essa verificação falha: participantes aceitam saídas erradas ou ignoram
desinformação (SB2-1255, RSL2-2710), editam mais o que aceitaram sem
verificar (RSL2-2593) e compreendem pior o código feito com IA (RSL2-0559,
SB2-0258). A revisão por IA aponta mais problemas sem achar mais defeitos
(RSL2-2484) e tem precisão baixa (SB2-0873). O quarto tema (TA4) repete o
contraste da Tabela~\ref{tab:direcao-medida-rsl2}: a percepção de qualidade
não coincide com a medição. Em RSL2-2650, a qualidade percebida foi maior com
IA, sem diferença medida. Em RSL2-0111, parte dos usuários percebe qualidade
maior, enquanto o indicador de depuração não mostra tendência. Muitos estudos
sustentam a direção atribuída só em autorrelato (RSL2-0826, RSL2-2903,
SB2-0846).

O quinto tema (TA5) é que contexto, iteração com retorno da execução e escopo
pequeno melhoram a saída. Há ganho medido com laço de correção guiado por
testes (RSL2-1320, RSL2-3904), com \textit{logs} como contexto (SB2-2369) e
com \textit{prompting} estruturado (SB2-1811). Sem contexto do projeto, as
saídas erram ou ficam genéricas (SB2-0162, RSL2-4151). O sexto tema (TA6) é
que a documentação gerada ajuda (RSL2-0613, SB2-0471), mas pede revisão de
precisão e de verbosidade (RSL2-1500, RSL2-2650, RSL2-1878).

Há contradições sobre manutenibilidade, com análise estática pior em
RSL2-1246, SB2-1605 e SB2-2107 e boa em RSL2-0487, RSL2-0488 e RSL2-1840, e
sobre tarefas complexas, em que SB2-2049 mede correção maior com IA e
RSL2-2839, RSL2-2068 e SB2-0279 medem queda. As lacunas principais são três.
Defeitos em produção e dívida técnica quase não são medidos (RSL2-0216,
SB2-0308, RSL2-0111). Poucos estudos acompanham a qualidade ao longo do tempo
(RSL2-0111, RSL2-0488, RSL2-4074). E a eficácia dos testes gerados em achar
defeitos raramente é medida (RSL2-2484, RSL2-0378).
% PENDENTE: o código C37 de B2 inclui RSL2-1593, cujo campo de qualidade só
% registra que "7/11 ignoraram o código" (caso limítrofe da conferência por
% script). R1 e R2 conferem.

% ---------------------------------------------------------------------------
\subsection{QP3 e QP4: segurança e custo}
\label{sec:resultados-rsl2-qp3-qp4}

% Gerado por docs/triagem/sintese/scripts/gera_tex.py. Não editar à mão.
\begin{table}[htbp]
\centering
\caption{Temas analíticos: RSL 2, QP3 e QP4}
\label{tab:temas-B3}
\small
\begin{tabular}{lp{11cm}r}
\hline
\textbf{Tema} & \textbf{Descrição} & \textbf{Estudos} \\
\hline
TA1 & Código gerado não é seguro por padrão, e o usuário não percebe & 39 \\
TA2 & Mitigações eficazes são ativas e em camadas; avisos passivos falham & 31 \\
TA3 & Confidencialidade dos dados é o risco que mais muda a governança & 41 \\
TA4 & Agentes ampliam a superfície de risco para o ambiente de desenvolvimento & 17 \\
TA5 & Custo é medido em unidades diferentes e quase nunca em dinheiro & 32 \\
TA6 & O custo oculto é o tempo humano de verificação & 24 \\
\hline
\end{tabular}
\fonte{Elaborada pelo autor. ``Estudos: número de estudos distintos nos códigos do tema.}
\end{table}


A Tabela~\ref{tab:temas-B3} reúne os temas de segurança (TA1 a TA4) e de
custo (TA5 e TA6). O primeiro tema (TA1) é que o código gerado não é seguro
por padrão, e o usuário não percebe. Nenhum estudo mediu ganho de segurança só
por usar IA. O único ganho apareceu com instrução explícita de segurança
(RSL2-0792). As medições divergem entre piora (RSL2-2576, RSL2-1246,
SB2-1605) e efeito nulo (RSL2-2812, SB2-2049, RSL2-0792), mas convergem num
ponto: quem usa IA superestima a segurança do próprio código (RSL2-2576,
SB2-2049, RSL2-1799). As falhas se concentram em classes conhecidas, como
injeção, criptografia, credenciais e dependências (RSL2-1328, RSL2-1246,
RSL2-0906). Com assistente envenenado, a adoção de código inseguro chega a
100\% na condição de referência (\textit{baseline}) de RSL2-1328. O segundo
tema (TA2) é que as mitigações que funcionam são ativas e em camadas, e os
avisos passivos falham. Em RSL2-1328, a interface Skeleton reduziu o uso
inseguro, e destaque e aviso não tiveram efeito significativo. A instrução
de segurança melhorou o resultado em RSL2-0792 e RSL2-1799, mas falhou com o
retorno do CodeQL em SB2-1605. A análise estática mede e filtra (RSL2-1246,
RSL2-0487, SB2-0308), mas deixou passar correções incorretas (SB2-0597). A
revisão humana é a mitigação mais recomendada (RSL2-0201, RSL2-2477,
SB2-0308), mas seu efeito não foi medido isoladamente.

O terceiro tema (TA3) é que a confidencialidade dos dados é o risco que mais
muda a governança. O risco mais citado é a exposição de código e dados ao
provedor (RSL2-0205, RSL2-2566, SB2-0678, SB2-1783). O único experimento
sobre isso mediu mais compartilhamento de dados sensíveis com ChatGPT
(RSL2-1182). As organizações reagem com proibições e listas aprovadas
(RSL2-0225, SB2-0873), modelos internos (RSL2-0065, RSL2-2477, SB2-0308) e
anonimização (RSL2-0246, RSL2-1546), sem avaliação quantitativa da
anonimização. Proibir sem oferecer alternativa leva ao uso escondido
(RSL2-0225, SB2-1783). O quarto tema (TA4) é que agentes ampliam a superfície
de risco para o ambiente de desenvolvimento. O risco passa do trecho de
código para ações no ambiente: remoção de testes e reversão de mudanças
(RSL2-0150), comandos não pedidos que apagam \textit{logs} (RSL2-2485) e
alterações não intencionais (RSL2-0750, SB2-0597). As mitigações relatadas
são de engenharia do ambiente: \textit{sandbox} e validação determinística
(RSL2-0150), menor privilégio e nenhum \textit{merge} automático (RSL2-3381).

Quanto ao custo, o quinto tema (TA5) é que ele é medido em unidades
diferentes e quase nunca em dinheiro. Os estudos usam \textit{tokens}
(RSL2-0150, SB2-1193, SB2-2374), dólares (RSL2-0378, RSL2-0487, RSL2-3833),
\textit{hardware} (RSL2-0729, SB2-0206), tempo de máquina (SB2-0844) e
iterações (SB2-0162). Um agente genérico usou 12,11 vezes mais
\textit{tokens} que uma ferramenta dedicada (SB2-2374), e o mesmo
\textit{pipeline} custou US\$~0,24 com Devstral e US\$~6,17 com GPT-o1
(RSL2-0378). A redução de 96 a 98\% relatada em RSL2-0487 compara o custo
real com uma estimativa, e não com uma medida. O sexto tema (TA6) é que o
custo oculto é o tempo humano de verificação (RSL2-0240, RSL2-0214,
RSL2-2484). Nas organizações, licenças e retorno incerto são barreiras
percebidas (RSL2-1668, SB2-1783, RSL2-2624).

As contradições tratam do efeito medido sobre a segurança, que piora em
alguns experimentos e não muda em outros, e do custo, visto como barreira em
alguns estudos (RSL2-1668, SB2-1783) e como compensado pelos ganhos em outro,
sem medida (RSL2-2477). As lacunas são claras. Nenhum estudo mediu consumo de
energia, que só aparece como menção (RSL2-3733, SB2-0873). Nenhum mediu
vulnerabilidades introduzidas por agentes autônomos. E vários experimentos de segurança usam assistentes envenenados de propósito ou
modelos antigos (RSL2-1328, RSL2-2843, RSL2-1799, RSL2-2576), o que limita a
generalização para as ferramentas comerciais atuais.

% ---------------------------------------------------------------------------
\subsection{QP5: processo e governança}
\label{sec:resultados-rsl2-qp5}

% Gerado por docs/triagem/sintese/scripts/gera_tex.py. Não editar à mão.
\begin{table}[htbp]
\centering
\caption{Temas analíticos: RSL 2, QP5 e QP6}
\label{tab:temas-B4}
\small
\begin{tabular}{lp{11cm}r}
\hline
\textbf{Tema} & \textbf{Descrição} & \textbf{Estudos} \\
\hline
TA1 & O portão humano permanece; a verificação muda de lugar & 129 \\
TA2 & A verificação é o novo gargalo e falha quando é só declarada & 97 \\
TA3 & Governança explícita, escrita e rastreável & 43 \\
TA4 & O desenvolvedor vira orquestrador, com risco de perder competência e interlocução & 73 \\
TA5 & Segurança não surge sozinha: precisa ser pedida e verificada & 87 \\
TA6 & A evidência sobre ferramentas e autonomia é heterogênea, datada e pouco transferível & 129 \\
\hline
\end{tabular}
\fonte{Elaborada pelo autor. ``Estudos: número de estudos distintos nos códigos do tema.}
\end{table}


A leitura B4 cobre a QP5 e a QP6 (Tabela~\ref{tab:temas-B4}). Os temas TA1 e
TA6 somam os 129 estudos porque incluem os códigos de mapeamento (ferramenta,
autonomia e contexto), que todo estudo recebe. Esses números indicam
cobertura, e não apoio a cada afirmação. O mesmo vale para TA5, cujo total
vem em grande parte do código de revisão humana, e não dos estudos que tratam
de segurança no processo.

O primeiro tema (TA1) é que o portão humano permanece e a verificação muda de
lugar. Nos agentes implantados, a extração sempre descreve uma decisão humana
antes da integração (RSL2-0150, SB2-0162, SB2-2613, SB2-1193), e os
desenvolvedores não querem aprovação nem \textit{commit} automáticos
(RSL2-3381). Testes, análise estática, juiz LLM e revisão cruzada entre
modelos passam a filtrar antes do humano (SB2-0597, RSL2-0383, RSL2-0487,
SB2-0308). O segundo tema (TA2) é que a verificação vira o gargalo e falha
quando é só declarada. O esforço sai da escrita e vai para verificar,
integrar e reescrever \textit{prompts}, o que foi medido por telemetria
(RSL2-0111), observado em experimento (RSL2-0240) e percebido em
\textit{survey} (RSL2-4074). O comportamento observado contradiz o declarado:
em RSL2-1799, 25 de 33 participantes disseram verificar o código, mas vários
usaram MD5; outros aceitaram sugestões sem editar (RSL2-2216). Os estudos
respondem com tarefas decompostas (RSL2-0487, RSL2-3990, RSL2-3768),
delegação conforme a criticidade (RSL2-0761, SB2-1960, RSL2-4104) e tempo de
revisão reservado na estimativa (RSL2-0240).

O terceiro tema (TA3) é que a governança precisa ser escrita e rastreável.
Políticas de uso são raras ou tardias (RSL2-0065, RSL2-2903, RSL2-2477).
RSL2-0201 recomenda políticas internas, um roteiro de adoção e o alinhamento
à ISO/IEC 42001. A rastreabilidade aparece em práticas concretas:
\textit{prompts} versionados (RSL2-0488, RSL2-1546, RSL2-1186), registro de
\textit{prompts} e saídas (RSL2-0201), autodeclaração do código gerado
(RSL2-0657) e trajetória do agente no \textit{diff} (SB2-0162). Mas só 83
(19\%) das organizações que permitem IA generativa têm um processo de
documentação conhecido (SB2-0678). A responsabilidade continua humana
(RSL2-2477, SB2-1143). O quarto tema (TA4) é que o desenvolvedor vira
orquestrador, com risco de perder competência e interlocução. A mudança de
papel é sobretudo percebida (SB2-1143, RSL2-4074, SB2-0207). Os riscos de
dependência são observados em experimentos: dependência crescente ao longo
da tarefa (SB2-1818), menor processamento de informação (SB2-2163) e perda
de posse do código (RSL2-0559). O quinto tema (TA5) é que a segurança precisa
ser pedida e verificada. Desenvolvedores revisam só a funcionalidade
(RSL2-0792, RSL2-2843), e em RSL2-2843 91,9\% desconheciam ataques
de envenenamento. Aparecem no processo \textit{checklists} de
vulnerabilidade (RSL2-0906) e \textit{pipeline} com Fortify e revisão humana
(SB2-0308).

A principal contradição é a ordem da revisão. Vários estudos põem o LLM como
primeira passada (RSL2-0322, RSL2-4151, SB2-0132, SB2-2369), mas RSL2-2484 o
recomenda depois da revisão manual, porque ele ancora a atenção do revisor. A lacuna mais relevante
para este trabalho é que as práticas de governança (\textit{prompts}
versionados, rastreabilidade, autodeclaração e políticas) são propostas ou
descritas, mas nenhum estudo mede o efeito delas. Também não há estudo sobre
o desenvolvimento individual de um produto completo em contexto acadêmico,
como num TCC. Os casos mais próximos, RSL2-0487 e SB2-1960, são casos únicos.
Por fim, nenhum estudo mede a perda de habilidades ao longo do tempo em
desenvolvimento real.

% ---------------------------------------------------------------------------
\subsection{Práticas candidatas para o desenvolvimento do protótipo}
\label{sec:resultados-rsl2-praticas}

A síntese termina na rastreabilidade para o trabalho
(Seção~\ref{sec:metodologia-rsl}): os temas de B1 a B4 viram práticas
candidatas para o processo de desenvolvimento do buscador. São candidatas, e
não práticas comprovadas, porque nenhum estudo mediu o efeito das práticas de
governança e nenhum tratou do desenvolvimento individual num TCC. A lista
abaixo dá cada prática com os estudos que a apoiam.

\begin{enumerate}
	\item \textbf{Medir tempo e retrabalho no diário do desenvolvimento.}
	Registrar o tempo por tarefa, os descartes e o retrabalho, e não confiar
	na sensação de rapidez (B1, TA1 e TA3; B4, TA2): RSL2-0205, RSL2-0240,
	RSL2-0246, RSL2-2593, SB2-1818, RSL2-4034.
	\item \textbf{Testes, \textit{lint} e análise estática de segurança
	(SAST) como portão na integração contínua,} antes da revisão humana (B2,
	TA2; B3, TA2; B4, TA1 e TA5): RSL2-0487, SB2-0597, SB2-0308, RSL2-1186,
	RSL2-2566, SB2-0162, RSL2-1328, SB2-2049, RSL2-0383.
	\item \textbf{Testes escritos ou validados pelo autor como critério de
	aceitação} do código gerado (B4, TA1): RSL2-0180, RSL2-0906, SB2-1605,
	SB2-2107, RSL2-3904, SB2-1783.
	\item \textbf{Revisão humana de cada \textit{diff}.} O agente não faz
	\textit{commit}, \textit{merge} nem \textit{push} sem revisão, e a
	revisão por outro modelo complementa a humana sem substituí-la (B2, TA3;
	B4, TA1): RSL2-0150, RSL2-3381, SB2-0162, SB2-2613, RSL2-2485, RSL2-0750,
	SB2-2374, RSL2-0322, RSL2-4151, SB2-0132, RSL2-2484, SB2-0873.
	\item \textbf{Tarefas pequenas, com contexto e com execução no laço.}
	Decompor cada funcionalidade em itens revisáveis, dar ao assistente o
	contexto do repositório e validar por execução (B2, TA5; B4, TA2):
	RSL2-0487, RSL2-3990, RSL2-3768, RSL2-1320, RSL2-3904, SB2-2369,
	RSL2-2746, RSL2-2815.
	\item \textbf{Requisitos de segurança explícitos no \textit{prompt} e
	\textit{checklist} de revisão} para o código gerado (B3, TA1; B4, TA5):
	RSL2-0792, SB2-2049, RSL2-2576, RSL2-0906, RSL2-2843.
	\item \textbf{Nenhum dado real ou sigiloso nos \textit{prompts}.} Usar
	dados sintéticos, já que o buscador trata dados de interesse policial
	(B3, TA3; B4, TA3): RSL2-1182, SB2-1818, RSL2-4104, RSL2-1385, SB2-1143,
	RSL2-0205, RSL2-2566.
	\item \textbf{Agente em \textit{sandbox}.} Contêiner sem credenciais
	reais, menor privilégio, \textit{commits} pequenos e testes protegidos
	contra remoção (B3, TA4): RSL2-0150, RSL2-2485, RSL2-3381, RSL2-0750.
	\item \textbf{Arquitetura definida pelo autor e trechos feitos sem IA}
	(B4, TA2 e TA4): RSL2-0487, SB2-1960, RSL2-3381, RSL2-0393, RSL2-0180,
	RSL2-1385, RSL2-0559.
	\item \textbf{Política escrita de uso de IA e registro rastreável.}
	Registrar ferramenta, modelo, versão, data, nível de autonomia,
	\textit{prompts} e \textit{tokens} de cada sessão, e declarar no
	\textit{commit} o código gerado (B3, TA5; B4, TA3 e TA6): RSL2-0201,
	RSL2-0488, RSL2-1546, RSL2-1186, RSL2-0657, RSL2-0750, SB2-0933,
	SB2-1193, SB2-2374.
\end{enumerate}
% PENDENTE: R1 e R2 decidem quais práticas entram no processo do protótipo.
% Os IDs das práticas 1 e 8 vêm das sínteses dos temas citados (B1 e B3 não
% listam estudos nas implicações).
